% ============================================================
% §6 · Ejecución del software — Manual de Usuario K-QGMIP
% ============================================================

\section{Ejecución del software}

Esta sección describe en detalle cómo ejecutar K-QGMIP. Se organiza en tres partes:
la entrada de datos del sistema (común a ambos métodos), la ejecución del método KGeoMIP
y la ejecución del método KQNodes. Antes de proceder, asegúrese de haber completado
la instalación descrita en la Sección de Instalación y Configuración.

% ============================================================
\subsection{Entrada de datos del sistema}
\label{sec:entrada_datos}
% ============================================================

Ambos métodos requieren los mismos datos de entrada: una Matriz de Probabilidad de
Transición (TPM), un estado inicial y tres máscaras binarias que definen el subsistema
de interés.

% ------------------------------------------------------------
\subsubsection{Matriz de Probabilidad de Transición (TPM)}
% ------------------------------------------------------------

La \textbf{TPM} codifica la dinámica del sistema: para cada estado $s(t)$ del sistema en
el instante $t$, la TPM especifica la distribución de probabilidad sobre los estados
posibles en $t+1$. K-QGMIP utiliza la representación \textit{estado-nodo} en notación
\textit{little-endian}.

Las redes de prueba se distribuyen como archivos CSV con el esquema de nombre:

\begin{lstlisting}
N{n}{pagina}.csv
\end{lstlisting}

donde \texttt{n} es el número de nodos y \texttt{pagina} es una letra que distingue
distintas redes del mismo tamaño (\texttt{A}, \texttt{B}, etc.).
Por ejemplo, \texttt{N20A.csv} es la red de muestra de 20 nodos, página A.

La ubicación de los archivos depende del método utilizado:

\begin{table}[h!]
\centering
\begin{tabular}{|l|l|}
\hline
\textbf{Método} & \textbf{Carpeta de muestras} \\
\hline
KQNodes  & \texttt{code/QNodes/src/.samples/} \\
KGeoMIP  & \texttt{code/GeoMIP/data/samples/} \\
\hline
\end{tabular}
\caption{Ubicación de las redes de prueba (TPM en CSV) por método.}
\label{tab:rutas_samples}
\end{table}

En KGeoMIP, el \texttt{Manager} busca la red en \texttt{data/samples/} de forma
predeterminada; la ruta puede sobreescribirse con la variable de entorno
\texttt{GEOMIP\_SAMPLES\_DIR} si se desea apuntar a un directorio distinto.

La TPM se carga automáticamente a través de la clase \texttt{Manager}:

\begin{lstlisting}[language=Python]
import numpy as np
from src.controllers.manager import Manager

gestor = Manager("10000000000000000000")  # estado inicial n=20
tpm = gestor.cargar_red()  # np.ndarray de shape (2^n, n)
\end{lstlisting}

Internamente, \texttt{cargar\_red()} usa \texttt{numpy.genfromtxt} con delimitador
\texttt{;} y selecciona el archivo \texttt{N\{n\}\{pagina\}.csv} correspondiente al
estado inicial configurado.

% ------------------------------------------------------------
\subsubsection{Estado inicial del sistema}
% ------------------------------------------------------------

El \textbf{estado inicial} $s(t)$ es una cadena binaria de longitud $n$ que indica
el estado de cada nodo en el instante de análisis. El nodo más a la izquierda
corresponde al bit de menor peso (\textit{little-endian}).

\begin{lstlisting}[language=Python]
ESTADO = "10000000000000000000"   # n=20: nodo 0 activo, resto inactivos
ESTADO = "11111111"               # n=8:  todos los nodos activos
ESTADO = "101"                    # n=3:  nodos 0 y 2 activos
\end{lstlisting}

La longitud de la cadena determina automáticamente el tamaño $n$ del sistema.

% ------------------------------------------------------------
\subsubsection{Máscaras binarias}
% ------------------------------------------------------------

Tres cadenas binarias de longitud $n$ definen el subsistema que se analizará:

\begin{itemize}
  \item \textbf{Condición}: nodos cuyo estado en $t$ se fija al analizar el subsistema.
        Habitualmente se establece como \texttt{"1"*n} (todos los nodos).
  \item \textbf{Alcance} (\textit{scope}): subconjunto de nodos en $t+1$ que se incluyen
        en el análisis. Ejemplos para $n=5$: \texttt{"11100"} incluye los nodos A, B y C.
  \item \textbf{Mecanismo}: subconjunto de nodos en $t$ que actúan como causa del
        subsistema. Ejemplos para $n=5$: \texttt{"00111"} incluye los nodos C, D y E.
\end{itemize}

En los archivos de prueba por lotes (Excel), alcance y mecanismo se especifican
con letras (\texttt{"ABC"}, \texttt{"CDE"}) y el \textit{framework} los convierte
automáticamente a cadenas binarias.

% ------------------------------------------------------------
\subsubsection{Inicialización y preparación del subsistema}
% ------------------------------------------------------------

El objeto \texttt{Manager} centraliza la carga de la TPM y la configuración de rutas.
Una vez creado, se pasa la TPM directamente al constructor de la estrategia elegida:

\begin{lstlisting}[language=Python]
from src.controllers.manager import Manager
from src.strategies.kqnodes import KQNodes

gestor = Manager(ESTADO)
tpm    = gestor.cargar_red()
kqn    = KQNodes(tpm)
\end{lstlisting}

La preparación del subsistema (condicionamiento y substracción) es automática:
ocurre dentro del método \texttt{aplicar\_estrategia()} de la clase base \texttt{SIA}
y el usuario no interviene directamente en ese proceso.

% ============================================================
\subsection{Método KGeoMIP (Estrategia Geométrica)}
\label{sec:kgeomip}
% ============================================================

KGeoMIP extiende la estrategia geométrica de GeoMIP al caso de k-particiones mediante
un \textbf{refinamiento divisivo top-down} (variante E4): parte de la bipartición óptima
que GeoMIP ya encontró ($k=2$) y en cada paso divide iterativamente la parte cuyo corte
produce el menor incremento de pérdida $\Delta\Phi$, hasta obtener $k$ partes.
Toda la tabla de costos $T$ y la matriz de similitud causal $S$ se calculan \textbf{una
sola vez} y se reutilizan para todos los valores de $k$, lo que hace al método eficiente
incluso para sistemas de hasta 25 nodos.

% ------------------------------------------------------------
\subsubsection{Parámetros de configuración}
% ------------------------------------------------------------

Todos los parámetros se configuran en el archivo
\texttt{code/GeoMIP/exec\_kgeomip.py}:

\begin{lstlisting}[language=Python]
ESTADO:   str = "1" + "0" * 24   # estado inicial (define n=25)
K:        int = 5                 # numero de partes (2, 3, 4 o 5)
VARIANTE: str = "E4"             # heuristica de refinamiento
MUESTRA:  str = "A"              # pagina de la red de prueba
\end{lstlisting}

\begin{table}[H]
\centering
\begin{tabular}{|l|l|p{7cm}|}
\hline
\textbf{Parámetro} & \textbf{Valores válidos} & \textbf{Descripción} \\
\hline
\texttt{ESTADO} & Cadena binaria longitud $n$ &
  Estado inicial $s(t)$. Su longitud determina el tamaño del sistema y
  la red a cargar (\texttt{N\{n\}\{MUESTRA\}.csv}). \\
\texttt{K} & \texttt{2}, \texttt{3}, \texttt{4}, \texttt{5} &
  Número de partes de la k-partición. Con \texttt{K=2} reproduce
  exactamente el resultado de GeoMIP. \\
\texttt{VARIANTE} & \texttt{"E4"}, \texttt{"A"} &
  \texttt{"E4"} = refinamiento divisivo por MinHeap con $\Delta\Phi$
  (recomendado; garantiza regresión $k=2$ y monotonicidad). \newline
  \texttt{"A"} = clustering aglomerativo sobre la matriz de similitud $S$
  (solo para comparación experimental, sin garantías). \\
\texttt{MUESTRA} & \texttt{"A"}, \texttt{"B"}, \ldots &
  Letra que identifica la red de prueba. \\
\hline
\end{tabular}
\caption{Parámetros configurables de KGeoMIP en \texttt{exec\_kgeomip.py}.}
\label{tab:params_kgeomip}
\end{table}

% ------------------------------------------------------------
\subsubsection{Pasos de ejecución}
% ------------------------------------------------------------

\begin{enumerate}
  \item \textbf{Abrir el archivo de configuración.} Con el entorno virtual activo,
        navegue a la carpeta del método:
\begin{lstlisting}[language=bash]
cd code/GeoMIP
\end{lstlisting}

  \item \textbf{Configurar los parámetros.} Edite \texttt{ESTADO}, \texttt{K},
        \texttt{VARIANTE} y \texttt{MUESTRA} en \texttt{exec\_kgeomip.py} según
        el sistema y el análisis deseado.

  \item \textbf{Ejecutar el script.}
\begin{lstlisting}[language=bash]
uv run exec_kgeomip.py
\end{lstlisting}
        El \textit{framework} carga la TPM desde \texttt{data/samples/N\{n\}\{MUESTRA\}.csv},
        lee todas las filas de alcance y mecanismo del Excel de pruebas para el tamaño $n$,
        y ejecuta KGeoMIP sobre cada combinación con los parámetros configurados.

  \item \textbf{Localizar los resultados.} Los archivos de salida se guardan en:
\begin{lstlisting}
results/kgeomip/resultado__N{n}_{MUESTRA}_{K}.csv
results/kgeomip/resultado__N{n}_{MUESTRA}_{K}.md
\end{lstlisting}
        El CSV contiene una fila por cada combinación (prueba, $k$, variante) con las
        columnas: \texttt{Prueba}, \texttt{Alcance}, \texttt{Mecanismo}, \texttt{k},
        \texttt{Criterio}, \texttt{Partición}, \texttt{Pérdida ($\varphi$)},
        \texttt{Tiempo (s)}.
\end{enumerate}

% ------------------------------------------------------------
\subsubsection{Ejemplo de ejecución}
% ------------------------------------------------------------

El siguiente ejemplo ejecuta KGeoMIP sobre la red de 25 nodos (muestra A),
buscando la 5-partición de mínima información con la variante E4:

\begin{lstlisting}[language=Python]
# exec_kgeomip.py
ESTADO   = "1" + "0" * 24   # n=25, nodo 0 activo
K        = 5
VARIANTE = "E4"
MUESTRA  = "A"
\end{lstlisting}

Salida esperada en consola:

\begin{lstlisting}
n=25: 45 pruebas x variante E4
  Prueba   1/ 45 -- Alc: ABCDE    Mec: ABCDE
  Prueba   2/ 45 -- Alc: ABCDEF   Mec: ABCDE
  ...
  CSV: results/kgeomip/resultado__N25_A_5.csv (45 filas)
  MD:  results/kgeomip/resultado__N25_A_5.md
\end{lstlisting}

% ============================================================
\subsection{Método KQNodes (Estrategia Submodular)}
\label{sec:kqnodes_ejecucion}
% ============================================================

KQNodes extiende el algoritmo de Queyranne (Ordenamiento de Máxima Adyacencia)
al caso de k-particiones mediante refinamiento iterativo de biparticiones.
Es la estrategia prioritaria para sistemas de tamaño moderado a grande donde
la eficiencia computacional es crítica.

% ------------------------------------------------------------
\subsubsection{Parámetros de configuración}
% ------------------------------------------------------------

Todos los parámetros se configuran directamente en el archivo
\texttt{code/QNodes/exec\_kqnodes.py}:

\begin{lstlisting}[language=Python]
ESTADO:   str = "1" + "0" * 19   # estado inicial (define n=20)
K:        int = 5                 # numero de partes (2, 3, 4 o 5)
CRITERIO: str = "C4"             # criterio de refinamiento
MUESTRA:  str = "A"              # pagina de la red de prueba
\end{lstlisting}

\begin{table}[h!]
\centering
\begin{tabular}{|l|l|p{7cm}|}
\hline
\textbf{Parámetro} & \textbf{Valores válidos} & \textbf{Descripción} \\
\hline
\texttt{ESTADO}   & Cadena binaria de longitud $n$ & Estado inicial $s(t)$. Su longitud determina el tamaño del sistema. \\
\texttt{K}        & \texttt{2}, \texttt{3}, \texttt{4}, \texttt{5} & Número de partes de la k-partición. \\
\texttt{CRITERIO} & \texttt{"C4"}, \texttt{"C1"} & Criterio de refinamiento: \texttt{C4} = corte marginal mínimo; \texttt{C1} = tamaño máximo de parte. \\
\texttt{MUESTRA}  & \texttt{"A"}, \texttt{"B"}, \ldots & Letra de la red de prueba a analizar (\texttt{N\{n\}\{MUESTRA\}.csv}). \\
\hline
\end{tabular}
\caption{Parámetros configurables de KQNodes en \texttt{exec\_kqnodes.py}.}
\label{tab:params_kqnodes}
\end{table}

% ------------------------------------------------------------
\subsubsection{Pasos de ejecución}
% ------------------------------------------------------------

\begin{enumerate}
  \item \textbf{Abrir el archivo de configuración.} En un terminal con el entorno
        virtual activo, navegue a la carpeta del método:
\begin{lstlisting}[language=bash]
cd code/QNodes
\end{lstlisting}

  \item \textbf{Configurar los parámetros.} Edite las cuatro variables al inicio de
        \texttt{exec\_kqnodes.py} según el sistema y el análisis que desea realizar.

  \item \textbf{Ejecutar el script.}
\begin{lstlisting}[language=bash]
uv run exec_kqnodes.py
\end{lstlisting}
        El \textit{framework} carga automáticamente la TPM correspondiente desde
        \texttt{data/DatosPruebas2026\_1.xlsx}, procesa todas las filas de alcance y
        mecanismo de la hoja asociada al tamaño $n$, y ejecuta KQNodes para cada
        combinación con el $k$ y criterio configurados.

  \item \textbf{Localizar los resultados.} Los archivos de salida se guardan en:
\begin{lstlisting}
results/kqnodes/resultado__N{n}_{MUESTRA}_{K}.csv
results/kqnodes/resultado__N{n}_{MUESTRA}_{K}.md
\end{lstlisting}
        El CSV contiene una fila por cada combinación
        (prueba, $k$, criterio) con las columnas:
        \texttt{Prueba}, \texttt{Alcance}, \texttt{Mecanismo}, \texttt{k},
        \texttt{Criterio}, \texttt{Partición}, \texttt{Pérdida ($\varphi$)},
        \texttt{Tiempo (s)}.
\end{enumerate}

% ------------------------------------------------------------
\subsubsection{Ejemplo de ejecución}
% ------------------------------------------------------------

El siguiente ejemplo ejecuta KQNodes sobre la red de 20 nodos (muestra A),
buscando la 5-partición de mínima información con criterio de corte marginal mínimo:

\begin{lstlisting}[language=Python]
# exec_kqnodes.py
ESTADO   = "1" + "0" * 19   # n=20, nodo 0 activo
K        = 5
CRITERIO = "C4"
MUESTRA  = "A"
\end{lstlisting}

Salida esperada en consola:

\begin{lstlisting}
n=20: 45 pruebas x 1 k x 1 criterios = 45 ejecuciones
  Prueba   1/ 45 -- Alc: ABCDE    Mec: ABCDE
  Prueba   2/ 45 -- Alc: ABCDEF   Mec: ABCDE
  ...
  CSV: results/kqnodes/resultado__N20_A_5.csv (45 filas)
  MD:  results/kqnodes/resultado__N20_A_5.md
\end{lstlisting}
